\scanpage{35}
Considérons le schéma des \deleted{triplets} systèmes $\mathrm{RigKos}$ $(T,T',B,B')$ avec
\[
\begin{aligned}
&T,T' \text{ tores maximaux en opposition}\\
&B,B' \text{ des Borels contenant }T,T'\text{ respectivement.}
\end{aligned}
\qquad\big[(h,p)=1\big]
\]
\inserted{(\uline{Rigidifications de Kostant})}
C'est un schéma lisse sur $S$, sur lequel $G$ opère librement.
\[
\begin{aligned}
&\mathrm{RigKos}/G \simeq \text{schéma fini étale sur }S,\ \text{\deleted{d'ordre} de rang }{\unsure{w|W}}^2/\varphi(h)hz\\
&\hspace{2em}\left\{\begin{aligned}
&\text{où }w\text{ est l'ordre du gpe de Weyl}\\
&\text{\deleted{\unsure{et}} où }h\text{ est le \unsure{n$^{\mathrm b}$|n°} de Coxeter}\\
&\varphi\text{ l'indicatrice d'Euler}\\
&z\text{ ordre du centre}
\end{aligned}\right.
\end{aligned}
\]
Il est associé à une représentation du gpe $E$ des autom extérieurs dans un schéma \inserted{fini} étale sur $\operatorname{Spec}\mathbb Z[1/h]$ [le schéma des \uline{classes} de rigidifications de Kostant de $G$ sur $\operatorname{Spec}(\mathbb Z[1/h])$].
Le dernier schéma est \deleted{bel et bien}\inserted{il} \deleted{de} splitté : prenons un couple en \unsure{opposit|opposition},  défini par
\[
\big(\underbrace{\mu_h\mathrel{\overset{\kappa}{\text{---}}}T}_{\substack{\text{canonique}\\\text{défini par épinglage}}},\quad
D\text{ engendré par }\ e_{\alpha_1}\unsure{\,\text{---}\,|+\cdots+}e_{\alpha_r}+
\underbrace{e_{-\psi}}_{\substack{\text{défini au}\\\text{signe près}}}=e
\]
On trouve $T'$ comme centralisateur \deleted{\unsure{commun}} de cet élément.
On a \deleted{le gpe} Alors le schéma des Borels \inserted{$\unsure{Q|\text{il}}$} contenant $T$, isomorphe à $W_{\mathbb Z}$, et le schéma $Q'$ des Borels contenant $T'$, qui est très tordu, \deleted{\illegible{1}} c'est un torseur sous $W'=N(T')/T'$.
On a dans $W\times W'$ le sous groupe \inserted{$\Phi=N\cap N'$} extension de $(\mathbb Z/h\mathbb Z)^*$ par $\deleted{\unsure{\mu|\mu_h}}\,\unsure{\mathcal Z|Z|\mathfrak z}\times\mu_h$, où $\mathcal Z\subset W$ (centralisateur de $\kappa$) et $\mu_h\subset W'$. Il opère \inserted{librement} sur $Q\times Q'$ à travers $W\times W'$, et on forme $(Q\times Q')/\Phi$.
Le gpe des autom extérieurs \inserted{$E$} de $G$ y opère via son action par $\pm1$ sur $\mathfrak g_{-\psi}$, donc en fait, \deleted{\unsure{ou celle}} ou bien l'opération de $E$ est triviale, ou elle se fait par une autom can.
% archive sheet number: 34
